\section{Further Questions}

\subsection{Problems with Extending the Construction to the Spin Case}

The base spaces $M$ of the circle bundles $E \to M$ constructed in the previous Chapter are all totally non-spin manifolds. As we have seen in Chapter \ref{sec:psc} many constructions in positive scalar curvature geometry require the existence of a spin structure, so it may also be desirable to get examples of such circle bundles over spin manifolds. However, our methods do not straightforwardly extend to the spin case. Let us discuss the difficulties that arise. 

Suppose we are given an $S^1$-principal bundle $\pi \colon E \to M$ over a spin manifold $M$ (in particular $M$ is orientable). First of all we cannot conclude the existence of a \psc-metric on $E$ from Theorem \ref{psc:thm:existence_result_totally_non_spin} anymore, since the total space $E$ is itself a spin manifold. The reason is the following. There exists a splitting 
\[
TE \cong HE \oplus VE,
\]
of the tangent bundle $TE$ of $E$ into a \textit{vertical} bundle $VE$ and a \textit{horizontal} bundle $HE$. The bundle $VE$ is simply the kernel of the differential $d\pi \colon TE \to TM$ and we can then choose $HE$ as any complementary subbundle (for example by choosing a fiber metric on $TE$). By construction $d \pi \colon HE \to TM$ is an isomorphism on fibers, such that $HE \cong \pi^\ast TM$. In our case $VE$ is one-dimensional, such that $w_2(VE) = 0$ and therefore
\[
w_2(TE) = w_1(HE) \smile w_1(VE) + w_2(HE) = \pi^\ast w_1(TM) \smile w_1(VE) + \pi^\ast w_2(TM) = 0.
\]
So instead we are inclined to use the spin analogue of Theorem \ref{psc:thm:existence_result_totally_non_spin} to conclude the existence of a \psc-metric on $E$, namely:

\vspace{\baselineskip}
\noindent\textbf{Theorem \ref{psc:thm:existence_result_spin}.}\itshape\ 
Let $W^{n+1}$ be a cobordism from $M_0$ to $M_1$, $n \geq 5$, such that $W$ is spin. Suppose that $M_1$ is connected, the inlcusion $M_1 \to W$ induces an isomorphism on $\pi_1$ and that $M_0$ carries a psc-metric. Then $M_1$ also carries a psc-metric. In fact the psc-metric on $M_0$ can be extended to a psc-metric on $W$, which is of product form near the boundary. 
\normalfont
\vspace{\baselineskip}

\noindent Proceeding as before we would like to take the disk bundle $DE$ as a nullbordism of $E$ and show:
\begin{enumerate}[label=(\Roman*\ensuremath{^\prime})]
	\item \label{constr:property1_spin} $DE$ is spin,
	\item \label{constr:property2_spin} The inclusion $S^1 \to E$ is trivial on $\pi_1$.
\end{enumerate}
However, these properties contradict themselves. To see this we make the following observations. Let $L \to M$ be the complex line bundle associated to $E \to M$ and consider the inclusion $\iota \colon DE \to L$. Since the differential $d \iota \colon TDE \to TL$ is an isomorphism on fibers we conclude $TDE \cong \iota^\ast TL$. Thus 
\[
\iota^\ast w_2(TL) = w_2(TDE) = 0
\]
by property \ref{constr:property1_spin}. But $\iota$ is a homotopy equivalence so consequently $w_2(TL) = 0$. Let $p \colon L \to M$ be the bundle proejction. As before we have a splitting 
\[
TL \cong HL \oplus VL
\]
into a vertical and horizontal subbundle with $HL \cong p^\ast TM$. For vector bundles it is a standard fact that $VL = \mathrm{ker}(dp) \cong p^\ast E$, hence 
\[
0 = w_2(TL) = \pi^\ast w_2(HL) + \pi^\ast w_1(HL) \smile p^\ast w_1(L) + p^\ast w_2(L) = p^\ast w_2(L).
\]
Again we conclude $w_2(E) = 0$ since $p$ is a homotopy equivalence. Therefore the first chern class $c_1(L)$ lies in the kernel of the coefficient homomorphism
\[
H^2(M;\Z) \to H^2(M;\Z_2)
\]
by Proposition \ref{bundles:prop:connection_chern_stiefel_whitney_classes}. In other words, $c_1(L)$ can be written as $c_1(L) = 2 \alpha$ for some $\alpha \in H^2(M;\Z)$. To derive a contradiction from this we claim that the following diagram commutes up to sign
\[\begin{tikzcd}[column sep=6em,row sep=2.5em]
	{\pi_2(M)} & {\pi_1(S^1)} \\
	{H_2(M;\Z)} & {H_1(S^1;\Z)}
	\arrow["\partial", from=1-1, to=1-2]
	\arrow[from=1-1, to=2-1]
	\arrow["\cong", from=1-2, to=2-2]
	\arrow["{\langle c_1(L),-\rangle [S^1]}", from=2-1, to=2-2]
\end{tikzcd}\]
where $\partial$ is the boundary map from the long exact sequence of homotopy groups for the fibration $E \to M$ and the vertical maps are given by the Hurewicz homomorphism. Moreover $[S^1]$ denotes the fundamental class of $S^1$. Granting this claim for now, it then follows from $c_1(L) = 2\alpha$ that $\partial \colon \pi_2(M) \to \pi_1(S^1)$ cannot be surjective and consequently $\pi_1(S^1) \to \pi_1(E)$ is non-zero by exactness. Hence a bundle satisfying both \ref{constr:property1_spin} and \ref{constr:property2_spin} does not exist. To show commutativity of the above square let $\tau_1 \to \CP^\infty$ be the universal complex line bundle (compare Example \ref{bundles:ex:examples_for_classifiying_spaces} \ref{bundles:es:examples_for_classifying_spaces_special_case}) and let $f \colon M \to \CP^\infty$ be a classifiying map for $L \to M$. Then essentially commutativity is due to the fact that the corresponding square for the \emph{universal} principal bundle commutes and all maps are compatible with the classifiying map $f$. To make this precise we extend the diagram as follows 

\[\begin{tikzcd}[column sep=1.5em,row sep=large]
	{\pi_2(\CP^\infty)} &&&&& {\pi_1(S^1)} \\
	& {\pi_2(M)} &&& {\pi_1(S^1)} \\
	& {H_2(M;\Z)} &&& {H_1(S^1;\Z)} \\
	{H_2(\CP^\infty;\Z)} &&&&& {H_1(S^1;\Z)}
	\arrow["\partial", "\cong"', bend left = 6, from=1-1, to=1-6]
	\arrow["\cong"', bend right = 6, from=1-1, to=4-1]
	\arrow["\cong", bend left = 6, from=1-6, to=4-6]
	\arrow["{f_\ast}", from=2-2, to=1-1]
	\arrow["\partial", from=2-2, to=2-5]
	\arrow[from=2-2, to=3-2]
	\arrow["{\mathrm{id}}"', from=2-5, to=1-6]
	\arrow["\cong", from=2-5, to=3-5]
	\arrow["{\langle c_1(L),-\rangle [S^1]}", from=3-2, to=3-5]
	\arrow["{f_\ast}"', from=3-2, to=4-1]
	\arrow["{\mathrm{id}}", from=3-5, to=4-6]
	\arrow["{\langle c_1(\tau_1),-\rangle [S^1]}", bend right = 6, from=4-1, to=4-6]
\end{tikzcd}\]

Again the top horizontal map is the boundary map from the long exact sequence for the fibration $ES^1 \to \CP^\infty$ and the outer vertical maps are given by the Hurewicz homomorphism. Since $\CP^\infty$ is simply connected $\pi_2(\CP^\infty) \to H_2(\CP^\infty;\Z)$ is an isomorphism. Moreover, the boundary map $\partial \colon \pi_2(\CP^\infty) \to \pi_1(S^1)$ is an isomorphism by Example \ref{bundles:ex:examples_for_classifiying_spaces} \ref{bundles:ex:homotopy_groups_of_BG}. Clearly 1 commutes. Furthermore 2 commutes by naturality of the long exact sequence of homotopy groups and 3 commutes by naturality of the Hurewicz homomorphism. For commutativity of 4 we compute for $a \in H_2(M;\Z)$:
\[
\langle c_1(\tau_1),f_\ast a \rangle = \langle f^\ast c_1(\tau_1),a \rangle = \langle c_1(L),a \rangle,
\]
using the fact that $f$ is a classifiying map in the last equation. Finally we have to show commutativity of the outer square. From the universal coefficient theorem we obtain that 
\[
H^2(\CP^\infty;\Z) \to \mathrm{Hom}(H_2(\CP^\infty;\Z),\Z), \;\; \beta \mapsto \langle \beta, - \rangle 
\]
is an isomorphism. By definiton $c_1(\tau_1) \in H^2(\CP^\infty;\Z)$ is a generator, thus $\langle c_1(\tau_1), - \rangle $ is a generator of $\mathrm{Hom}(H_2(\CP^\infty;\Z),\Z)$ and consequently evaluates to $\pm 1$ on a generator of $H_2(\CP^\infty;\Z)$. From this we see that
\[
\langle c_1(\tau_1), - \rangle [S^1] \colon H_2(\CP^\infty;\Z) \to H_2(S^1;\Z)
\]
is an isomorphism. So all maps in the outer square are isomorphisms and, moreover, all groups are isomorphic to $\Z$. Hence the outer square commutes up to sign and the claim follows. 



\subsection{A Question about the Homotopy Groups of the Space of \psc-Metrics}